\documentclass{ctexart}
\usepackage{amsmath,booktabs,graphicx}
\begin{document}
\section{模型}
为回答 Q1，先定义状态量与时间范围。
\begin{equation}
  y_t = f(x_t;\theta)+\varepsilon_t,\qquad t\in\mathcal T
  \label{eq:q1-observation}
\end{equation}
其中 $y_t$ 为观测量，单位与数据字典一致。
\begin{align}
  \widehat{\theta} &= \arg\min_{\theta\in\Theta} L(\theta),
    \label{eq:q1-fit}\\
  L(\theta) &= \sum_{t\in\mathcal T} (y_t-f(x_t;\theta))^2.
    \label{eq:q1-loss}
\end{align}
由式~\eqref{eq:q1-fit} 得到参数估计，再用式~\eqref{eq:q1-loss} 计算误差。
\begin{table}[htbp]
  \centering
  \caption{Q1 参数与单位}
  \label{tab:q1-parameters}
  \begin{tabular}{@{}ll@{}}
    \toprule
    符号 & 单位 \\
    \midrule
    $y_t$ & 1 \\
    \bottomrule
  \end{tabular}
\end{table}
\begin{figure}[htbp]
  \centering
  \rule{0.5\linewidth}{0.2\linewidth}
  \caption{Q1 误差示意}
  \label{fig:q1-error}
\end{figure}
图~\ref{fig:q1-error} 与表~\ref{tab:q1-parameters} 用于说明结果范围。
\end{document}
